% Generated draft: requires mathematical and source review.
\documentclass[11pt]{article}
\usepackage[margin=25mm]{geometry}
\usepackage{fontspec}
\setmainfont{Latin Modern Roman}
\usepackage{amsmath,amssymb,mathtools,unicode-math}
\setmathfont{Latin Modern Math}
\usepackage{enumitem,needspace,xurl,hyperref}
\setlength{\parindent}{0pt}
\setlength{\parskip}{5pt}
\setlength{\emergencystretch}{3em}
\newcommand{\N}{\mathbb{N}}
\newcommand{\Z}{\mathbb{Z}}
\newcommand{\Q}{\mathbb{Q}}
\newcommand{\R}{\mathbb{R}}
\newcommand{\C}{\mathbb{C}}
\newcommand{\F}{\mathbb{F}}
\newcommand{\A}{\mathbb{A}}
\newcommand{\PP}{\mathbb{P}}
\newcommand{\E}{\mathbb{E}}
\newcommand{\eps}{\varepsilon}
\newcommand{\dd}{\,\mathrm{d}}
\begin{document}
\section*{UnsolvedMath ID: 801193}
\textbf{Source problem code:} catalogue entry
\subsection{Definitions and mathematical statement}
Let \(n\geq 1\), let \(v\) be an indeterminate, and put \(q=v^2\). The following formulation uses the conventional normalization described below; the supplied record does not explicitly specify its normalization conventions.

Let \(H_n(q)\) be the algebra over \(\Q(v)\) generated by \(T_1,\ldots,T_{n-1}\), with the type \(A_{n-1}\) braid relations and quadratic relations
\[
(T_i-q)(T_i+1)=0.
\]
For \(w\in S_n\), let \(\ell(w)\) be its number of inversions and let \(T_w\) be the product along a reduced expression for \(w\). Write \(x<w\) for strict Bruhat order, characterized by the reduced-subword condition. Set \(\widetilde T_w=v^{-\ell(w)}T_w\). The Kazhdan--Lusztig element \(C'_w\) is characterized by
\[
\overline{C'_w}=C'_w,
\qquad
C'_w\in \widetilde T_w+
\sum_{x<w}v^{-1}\Z[v^{-1}]\,\widetilde T_x,
\]
where the bar involution sends \(v\) to \(v^{-1}\) and \(T_w\) to \(T_{w^{-1}}^{-1}\).

For each partition \(\mu\vdash n\), let \(\chi_q^\mu\) be the generic irreducible character labeled by the corresponding Specht module. Define the characteristic isomorphism from their \(\Q(v)\)-linear span to the degree-\(n\) symmetric functions by \(\operatorname{ch}_q(\chi_q^\mu)=s_\mu\), where \(s_\mu\) is the Schur symmetric function. For \(\lambda\vdash n\), let \(m_\lambda\) be the monomial symmetric function, the sum of the distinct monomials whose nonzero exponent multiset is \(\lambda\). The monomial virtual character is
\[
\phi^\lambda=\operatorname{ch}_q^{-1}(m_\lambda).
\]
Thus it is an integral linear combination of irreducible characters, with coefficients that need not be nonnegative.

The conjecture attributed in the supplied record to Mark Haiman asks whether, for every \(n\geq1\), every \(\lambda\vdash n\), and every \(w\in S_n\), writing \(d=\ell(w)\), one has
\[
F_{\lambda,w}(q):=\phi^\lambda\!\left(v^d C'_w\right)
=\sum_{j=0}^{d}a_jq^j,
\qquad a_j\in\Z_{\geq0},
\]
with
\[
a_j=a_{d-j}\quad(0\leq j\leq d),
\qquad
 a_0\leq a_1\leq\cdots\leq a_{\lfloor d/2\rfloor}.
\]
Coefficients are padded with zeros through exponent \(d\); the zero polynomial is permitted. The symmetry condition is equivalently \(q^dF_{\lambda,w}(q^{-1})=F_{\lambda,w}(q)\), expressing the source's symmetry about \(q^{\ell(w)/2}\). The compatibility of the explicit Hecke-algebra and character conventions above with the cited source requires human verification. The registry labels the problem open; this is a supplied status, not an independently verified mathematical assessment.
\subsection{Short English statement}
For every partition \(\lambda\vdash n\) and permutation \(w\in S_n\), Haiman's conjecture asks whether evaluating the monomial virtual character on \(q^{\ell(w)/2}C'_w\) produces a polynomial with nonnegative integer coefficients, symmetric about exponent \(\ell(w)/2\) and unimodal. The explicit normalization requires checking against the source, and the supplied open status has not been independently verified.
\subsection{Sources}
\begin{itemize}
\item Local UnsolvedMath registry, permanent ID \(801193\), exact source problem code \texttt{catalogue entry}; supplied title: ``Stanley's Positivity Problems in Algebraic Combinatorics --- Problem 19: Haiman's unimodality and positivity conjecture for Hecke algebra monomial characters on Kazhdan--Lusztig basis elements.'' The registry lists Stanley as proposer, year 1999, and status open.
\item Stanley, \emph{Positivity Problems and Conjectures in Algebraic Combinatorics} (1999), Problem 19. Supplied source URL: \url{https://math.mit.edu/~rstan/papers/problems.pdf}. The supplied provenance records access on 2026-10-16 and extraction as \texttt{pdf+html\_verified}; these are supplied provenance claims, not checks performed for this draft.
\item Mark Haiman, \emph{J. Amer. Math. Soc.} \textbf{6} (1993), 569--595. The supplied notes attribute the conjecture to Haiman and identify it with Stanley's Problem 19. No article title or separate URL is supplied for this reference.
\end{itemize}
\end{document}
